%Chapter 1
\section{Introduction}
\label{sec:introduction}

%Chapter 1.1
\subsection{Context and Motivation}
\label{subsec:context-motivation}

Secure communication is part of everyday infrastructure, but quantum computing threatens the public-key assumptions on which secure messaging depends. Signal has begun migrating its protocol stack towards post-quantum cryptography. This chapter motivates a closer examination of its initial-handshake choices, states the research problem and questions, and defines the contributions, scope, and structure of the dissertation.

The threat in question is the prospect of a sufficiently powerful quantum computer. Peter Shor's 1994 algorithm \cite{shor_algorithms_1994} demonstrated that a quantum computer capable of running it at sufficient scale would solve the integer factorisation problem and the discrete logarithm problem in polynomial time. The public-key cryptosystems that protect the overwhelming majority of internet traffic, including RSA and elliptic-curve cryptography, derive their security from the assumed classical intractability of exactly these problems. A cryptographically relevant quantum computer, once built, would therefore render these cryptosystems insecure.

Whether such a machine will be built, and if so when, remains uncertain. Current engineering estimates place the threshold far enough in the future that transitional planning is possible but close enough that it is not optional \cite{mosca_cybersecurity_2018}. The concrete danger, however, does not wait for the machine itself. Under the harvest-now-decrypt-later (HNDL) threat model, an adversary intercepts and stores encrypted traffic today, with the intention of decrypting it later once a suitable quantum computer becomes available \cite{mosca_cybersecurity_2018}. Any communication whose confidentiality needs to outlive the gap between now and that future moment is therefore effectively exposed to quantum attack already, at the time it is transmitted, regardless of the protocol's strength against present-day adversaries.

Signal's response to HNDL began in 2023 with PQXDH (Post-Quantum Extended Triple Diffie-Hellman), which augments X3DH with CRYSTALS-Kyber, since standardised as ML-KEM \cite{kret_pqxdh_nodate,kret_quantum_nodate}. PQXDH combines its classical and post-quantum shared secrets through a key derivation function, retaining a classical fallback while protecting recorded handshakes against future quantum attack.

PQXDH has received independent formal analysis \cite{bhargavan_formal_2024}, but the public record does not comparatively characterise Signal's primitive choices within the wider post-quantum landscape. That unresolved choice motivates this dissertation.

%Chapter 1.2
\subsection{Problem Statement}
\label{subsec:problem-statement}

Signal's PQXDH protocol combines X25519 alongside ML-KEM-1024 in a hybrid construction \cite{kret_pqxdh_nodate}. Signal's published rationale motivates this hybrid approach as a conservative safety measure against crypto-analytic progress on either primitive and describes ML-KEM as resting on solid foundations \cite{kret_pqxdh_nodate, kret_quantum_nodate}. What Signal has not published, however, is any comparative evaluation of alternative post-quantum KEMs or signature schemes in its deployment context, nor a quantitative analysis of how its specific constraints (namely asynchronous messaging, mobile-first clients, and server-side pre-key bundle storage at scale) bear on this choice.

This is not a library-only comparison: PQXDH requires the KEM shared secret to be bound to its public-key context, so every candidate must be assessed against protocol-level as well as deployment-level requirements \cite{bhargavan_formal_2024}.

This dissertation addresses that gap. It implements the PQXDH protocol with alternative post-quantum primitive combinations and benchmarks their handshake latency, protocol-object sizes, and modelled server-held cryptographic payload against a pinned PQXDH baseline, upstream libsignal v0.73.3, whose handshake uses the Kyber-1024 parameter set. The alternative instantiations built for this work are fully post-quantum at the handshake bootstrap and its transcript authentication, with the X25519 component removed from the key-derivation input rather than augmented. The Double Ratchet that follows continues to mix an X25519 agreement into the chain protecting the first encrypted payload, so the claim is bounded at that handshake boundary throughout the dissertation. The justification for this choice, together with the criteria used to select candidate primitives, is developed in Chapter 3.

%Chapter 1.3
\subsection{Aim and Research Questions}
\label{subsec:aims-objectives}

The overarching aim of this project is to examine the design space of post-quantum primitive selection for Signal's PQXDH key agreement protocol \cite{kret_pqxdh_nodate}, both through structured analysis of the cryptographic literature and through empirical measurement of a working implementation, and to situate Signal's currently-deployed primitive choices within that design space on the basis of comparative evidence rather than inference.

%Chapter 1.3.2
\subsubsection{Research Questions}
\label{subsubsec:research-questions}

The work is organised around three research questions.

\textbf{RQ1} -- Which post-quantum primitive combinations constitute plausible candidates for instantiating PQXDH, and what criteria should govern their selection given Signal's deployment context?

\textbf{RQ2} -- How do the selected post-quantum primitive combinations compare, on Signal-relevant metrics, both to each other and to the pinned PQXDH baseline defined in Chapter 3 \cite{kret_pqxdh_nodate}?

\textbf{RQ3} -- What does the comparative evidence reveal about the design space of PQXDH, and what implications does it carry for the primitive choices made in the deployed protocol and for the wider class of post-quantum secure messaging protocols?

%Chapter 1.3.3
\subsubsection{Contributions}
\label{subsubsec:introduction-contributions}

This dissertation makes five contributions.

C1. A structured literature review of post-quantum cryptographic primitives and of how widely deployed cryptographic protocols, including TLS 1.3, SSH, MLS, and iMessage PQ3, have approached post-quantum migration, synthesising the design patterns, trade-offs, and lessons that those communities have documented and mapping them onto Signal's context.

C2. An analysis of Signal's published design rationale for PQXDH \cite{kret_pqxdh_nodate, kret_quantum_nodate} and a corresponding selection framework for evaluating alternative post-quantum primitives as candidate replacements, articulating the protocol-level requirements that PQXDH imposes, including those identified by the formal analysis of Bhargavan et al. \cite{bhargavan_formal_2024}, together with the deployment-specific criteria against which each candidate must be judged.

C3. A working implementation of the PQXDH protocol instantiated with alternative fully post-quantum primitive combinations, constructed by integrating published post-quantum cryptographic libraries with the structure of the PQXDH specification \cite{kret_pqxdh_nodate}.

C4. A comparative benchmark of the implemented instantiations against the pinned PQXDH baseline of Chapter 3, produced under a documented and reproducible experimental methodology, and reporting handshake latency, per-primitive operation cost, serialised protocol-object sizes, a serialised cryptographic bundle-component total, and a modelled server-held cryptographic payload per device.

C5. An analysis of the resulting comparative design space that positions Signal's currently-deployed primitive choices within the observed performance and cost landscape, identifies the conditions under which alternative choices would be preferable, and draws implications for the wider class of post-quantum secure messaging protocols.

Contributions C1 and C2 primarily address RQ1, contributions C3 and C4 primarily address RQ2, and contribution C5 primarily addresses RQ3.

%Chapter 1.4
\subsection{Scope and Limitations}
\label{subsec:scope-limitations}

%Chapter 1.4.1
\subsubsection{Scope}
\label{subsubsec:scope}

The scope of this project is defined by three boundaries.

First, the implementation artefact addresses only the PQXDH key agreement protocol \cite{kret_pqxdh_nodate}. The Double Ratchet protocol, the Sender Keys protocol for group messaging, the Safety Numbers protocol for key verification, and the encrypted voice and video call protocol all fall outside the scope of the artefact, and no modifications to them are proposed or implemented. Signal's more recent Triple Ratchet and Sparse Post-Quantum Ratchet (SPQR) designs, which address post-compromise security in the continuous ratchet rather than in the initial handshake, are likewise outside the artefact's scope. The literature review in Chapter 2 does, however, discuss these other protocols and designs, both to situate PQXDH within Signal's broader post-quantum roadmap and to extract lessons from how the wider cryptographic community has approached post-quantum migration in other deployed protocols.

Second, the artefact substitutes post-quantum primitives within the structure of the published PQXDH specification. The resulting variants are recognisably PQXDH; what changes between variants is the choice of post-quantum KEM used to generate the post-quantum shared secret, together with the domain-separation labels bound to it; both variants use ML-DSA-87 to sign pre-keys and the transcript, and both remove the classical X25519 component from the handshake secret. The project does not propose a new key agreement protocol and preserves PQXDH's asynchronous workflow, in which the initiator fetches a published pre-key bundle and establishes a session with a single message. It does change the authenticated fields of that message and the post-quantum pre-key policy, so the resulting variants are PQXDH-derived and do not interoperate with Signal's deployed implementation.

Third, the artefacts support empirical comparison rather than production deployment. The evaluation covers computation, bandwidth, and modelled storage but provides no new protocol proof; because authentication and KDF inputs change, existing PQXDH analyses do not transfer. Findings are therefore classified as demonstrated, conditional, unresolved, or unmet.

%Chapter 1.4.2
\subsubsection{Limitations}
\label{subsubsec:introduction-limitations}

The principal limitations are two consumer platforms rather than mobile hardware, library maturity confounded with primitive choice, and theoretically motivated rather than exhaustive candidate selection. They are examined in Sections~\ref{subsubsec:lit-primitives},~\ref{subsec:selection}, and~\ref{subsec:threats-validity}.

%Chapter 1.5
\subsection{Structure}
\label{subsec:structure}

The remainder of this dissertation is organised as follows.

Chapter 2 establishes the protocol background, compares candidate primitives and deployed migrations, and derives the research gap.

Chapter 3 defines the requirements and threat model, selects the primitives, and specifies and assesses the two variants.

Chapter 4 defines reproducible correctness, performance, size, and storage evaluation procedures.

Chapter 5 reports the evidence and evaluates the variants against the stated requirements.

Chapter 6 addresses the project's legal, social, ethical, and professional implications.

Chapter 7 answers the research questions, consolidates limitations, and identifies future work.
